\begin{definition} \label{definition:eigenvalue_and_eigenvector_of_an_endomorphism_of_a_vector_space_over_a_field}
Let $F$ be a \CrefAndHyperrefIfExist{definition:field}{field}, let $V$ be an \CrefAndHyperrefIfExist{definition:vector_space_over_a_field}{$F$-vector space}, and let $\phi : V \to V$ be an \CrefAndHyperrefIfExist{definition:morphism_of_vector_spaces}{$F$-linear endomorphism} of $V$.
An element $\lambda \in F$ is called an \hldef{eigenvalue} of $\phi$ if there exists a nonzero vector $v \in V$, $v \neq 0$, such that
$\phi(v) = \lambda v.$
Any nonzero vector $v \in V$ satisfying $\phi(v) = \lambda v$ for some eigenvalue $\lambda \in F$ is called an \hldef{eigenvector} of $\phi$ corresponding to (or associated with) the eigenvalue $\lambda$.
\end{definition}
